\section{The first section}
\lipsum[1-2]\cite{jirauschek2014}

\section{Another section, this time with a long title}
\lipsum[3-4]
\begin{figure*}[tbp]
  \centering
  \includegraphics[width=.8\textwidth, height=.3\textheight]{example-image-a}
  \caption{This is an example image that should be printed
    at the top of a page.}
  \label{fig:one}
\end{figure*}
\lipsum[5-7]

\subsection{Sections can have subsections}
\lipsum[8-10]

\subsection{Another subsection}
\lipsum[11]
\begin{equation}
  \label{eq:one}
  a^2 + b^2 = c^2
\end{equation}
\lipsum[12]

\section{Some nice theory}
\lipsum[13]\cite{feynman1958theory}

\subsection{A particular use-case}
\lipsum[14]
\begin{table}
  \centering
  \caption{This is a table caption at the top of a table.}
  \label{tab:one}
  \begin{tabular}{lcr}
    \toprule
    Header 1 & Header 2 & Header 3 \\
    \midrule
    Data 1   & 3.5      & 4.3      \\
    Data 2   & 3.5      & 4.3      \\
    Data 3   & 3.5      & 4.3      \\
    Data 4   & 3.5      & 4.3      \\
    Data 5   & 3.5      & 4.3      \\
    \bottomrule
  \end{tabular}
\end{table}
\lipsum[15-18]

\subsection{Problems with the theory}
\lipsum[19-25]

\section{A better theory}
\lipsum[26-28]\cite{bardeen1973four}
\begin{figure*}[tbp]
  \begin{subfigure}[t]{.4\textwidth}
    \centering
    \includegraphics[width=0.8\textwidth]{example-image-b}
    \subcaption{Subfigure left, with its own caption.}
    \label{fig:two_a}
  \end{subfigure}
  \hfill
  \begin{subfigure}[t]{.5\textwidth}
    \centering
    \includegraphics[width=0.8\textwidth]{example-image-c}
    \subcaption{Subfigure right, with its own caption
      \cite{einstein1905elektrodynamik}.}
    \label{fig:two_b}
  \end{subfigure}
  \centering
  \caption{Another caption for the entire figure. The figure consists of
    two subfigures.}
  \label{fig:two}
\end{figure*}

\subsection{A long subsection title that spans multiple lines is not a good
  idea but it is possible nevertheless}
\lipsum[29-30]
\begin{align}
  \nabla\cdot\vec{D}  & = \rho                                         \\
  \nabla\cdot\vec{B}  & = 0                                            \\
  \nabla\times\vec{E} & = -\frac{\partial\vec{B}}{\partial t}          \\
  \nabla\times\vec{H} & = \vec{J} + \frac{\partial\vec{D}}{\partial t}
\end{align}
\lipsum[31]
\begin{gather}
  \oiint_{\partial\Omega}\vec{D}\times\mathrm{d}\vec{S}
  = \iiint_\Omega\rho\mathrm{d}V                            \\
  \oiint_{\partial\Omega}\vec{B}\times\mathrm{d}\vec{S} = 0 \\
  \oint_{\partial\Sigma}\vec{E}\mathrm{d}l
  = -\frac{d}{dt}\iint_\Sigma\vec{B}\mathrm{d}\vec{S}       \\
  \oint_{\partial\Sigma}\vec{H}\mathrm{d}l
  = \iint_\Sigma\vec{J}\mathrm{d}\vec{S}
  +\frac{d}{dt}\iint_\Sigma\vec{D}\mathrm{d}\vec{S}
\end{gather}
\lipsum[32-35]

\subsection{The last subsection of this chapter}
\lipsum[36-40]
